\chapter{Integración de inteligencia artificial en la supervisión infantil móvil}
\label{chap:ia-supervision-movil}

\insertminitoc
\parindent1.5em

\section{Técnicas de detección de riesgos digitales y análisis multimodal}

La detección de riesgos en el entorno móvil contemporáneo exige un enfoque multimodal que combine el Procesamiento de Lenguaje Natural (\gls{nlp}) con la Visión Artificial (CV). A diferencia de las herramientas tradicionales basadas en listas negras de URLs o filtros DNS, que han demostrado ser ineficaces frente a aplicaciones de mensajería cifrada y navegadores alternativos \cite{rojas_technical_2025}, la integración de modelos de Inteligencia Artificial permite una inspección profunda del contenido semántico y visual. Esta transición es vital, ya que la literatura técnica sugiere que los sistemas más robustos son aquellos que analizan el contenido en su forma final en la capa de Interfaz de Usuario (HCI), evaluando conjuntamente el texto, las imágenes y el audio tal como se presentan ante el menor \cite{jevremovic_keeping_2021}. Este enfoque "en pantalla" permite interceptar riesgos dentro de aplicaciones cerradas que tradicionalmente bloquean la inspección de red.\\


La visión artificial juega un rol determinante en la detección de contenido inapropiado que se presenta de forma efímera o camuflada dentro de aplicaciones de redes sociales y plataformas de video. Investigaciones recientes sobre la seguridad infantil en plataformas como YouTube destacan la necesidad de utilizar aprendizaje profundo (\textit{deep learning}) para clasificar automáticamente videos que eluden los filtros básicos de edad \cite{alqahtani_childrens_2023}. Esta capacidad de análisis visual es crítica considerando que un porcentaje significativo de los menores se expone involuntariamente a contenido pornográfico o violento debido a algoritmos de recomendación mal calibrados \cite{yusuf_lets_2023}. Además, la medición del uso problemático de la tecnología en la infancia temprana requiere evaluar no solo el tiempo de pantalla, sino el tipo de estímulos visuales que el niño consume solitariamente \cite{selak_problematic_2025}.\\

En el ámbito del texto y las interacciones privadas, las aplicaciones de mensajería como Instagram Direct Messages (DMs) se han convertido en el vector principal de victimización. Con la inminente transición hacia el cifrado de extremo a extremo (E2E), la intercepción de mensajes en tránsito es inviable. Sin embargo, estudios han demostrado que la extracción de características basadas en metadatos y señales lingüísticas locales permite entrenar clasificadores de \textit{Machine Learning} que detectan conversaciones inseguras con alta precisión \cite{ali_getting_2023}. Utilizando redes neuronales convolucionales (\gls{cnn}) y modelos de \textit{Random Forest}, es posible identificar indicadores de riesgo sexual y acoso a nivel de conversación y de mensaje individual, superando el 88\% de precisión (AUC) \cite{razi_sliding_2023}. \\

Además, la detección inteligente debe ser capaz de evolucionar junto con el lenguaje y las dinámicas juveniles. Una revisión sistemática sobre enfoques computacionales centrados en el humano advierte que el monitoreo no debe ser meramente punitivo, sino contextualmente consciente, diferenciando entre la jerga coloquial inofensiva y las tácticas manipulativas de ingeniería social utilizadas en las fases iniciales del \textit{grooming} \cite{razi_human-centered_2021}. Esto se complementa con arquitecturas de minería de datos como \textit{ChildProtect}, que construyen diccionarios dinámicos para rastrear valencias hostiles y discursos de odio en la navegación web, alertando a los padres sobre signos de radicalización o violencia cibernética \cite{ameur_childprotect_2023}. La combinación de estos modelos proporciona una red de seguridad integral que comprende las complejidades del comportamiento digital adolescente \cite{page_psychological_2025}.\\

\section{Registro, gestión de evidencias y protección de la privacidad}

Una vez que el sistema de \gls{ia} detecta un riesgo potencial, la gestión de la evidencia digital se convierte en el pilar fundamental de la intervención parental. La "Bitácora de Evidencias" propuesta en este modelo no es un simple repositorio de datos, sino un registro criptográficamente protegido que debe cumplir con estrictos estándares de integridad. La seguridad en infraestructuras móviles modernas, especialmente ante la llegada de redes 6G y la explosión de información accesible, exige que cualquier dato sensible sea cifrado de extremo a extremo y procesado bajo mecanismos de privacidad diferencial \cite{nguyen_security_2021}. En nuestro modelo, el acceso restringido mediante un PIN secreto asegura que la evidencia sea revisada exclusivamente por el tutor legal, mitigando el riesgo de que el propio sistema se convierta en una vulnerabilidad.\\

Esta protección es una respuesta directa a los hallazgos críticos de la comunidad de ciberseguridad, que han denunciado cómo muchas herramientas comerciales de control parental actúan como "caballos de Troya", abusando de los permisos del sistema operativo para rastrear y comercializar la información de los menores \cite{feal_angel_2020}. Esta paradoja de privacidad no se limita a los teléfonos inteligentes, sino que se extiende a dispositivos del Internet de las Cosas (\gls{iot}) orientados a la infancia, como las aplicaciones de monitoreo de bebés, las cuales frecuentemente carecen de cifrado robusto y exponen transmisiones de video íntimas a redes públicas \cite{schmidt_assessing_2023}. Por consiguiente, el diseño de la bitácora debe garantizar que los datos nunca abandonen el dispositivo físico (\textit{on-device storage}), devolviendo la soberanía de la información a la familia.\\


Desde una perspectiva pedagógica, la gestión de estas evidencias debe proporcionar un marco de "insumos para el diálogo" en lugar de herramientas para el castigo. La efectividad de las aplicaciones de control parental aumenta drásticamente cuando su diseño soporta la comunicación transparente entre padres e hijos \cite{stoilova_parental_2024}. Investigaciones que analizan el espacio de diseño de estas aplicaciones revelan que características como la granularidad del control y la transparencia algorítmica influyen fuertemente en si el menor percibe la herramienta como protección o como castigo \cite{wang_protection_2021}. Al registrar únicamente "momentos clave", el sistema evita el monitoreo totalizador, fomentando la agencia del menor y ayudándole a internalizar habilidades de autorregulación \cite{dumaru_its_2024}.\\

Técnicamente, el almacenamiento de evidencias visuales en dispositivos móviles impone desafíos significativos de rendimiento y gestión de memoria. Para evitar la degradación del sistema, la aplicación debe implementar políticas de retención de datos inteligentes y compresión eficiente. Estudios sobre el rendimiento del \textit{kernel} de Android demuestran que la gestión inadecuada de procesos en segundo plano provoca latencias severas en el cambio de aplicaciones (\textit{app-switching}). Rectificar el reclamo de páginas del núcleo y balancear el uso de la memoria virtual (zRAM) permite que el servicio de monitoreo capture y almacene capturas de pantalla sin causar "thrashing" ni ralentizar el dispositivo del menor, asegurando que la supervisión sea continua e imperceptible \cite{chou_rectifying_2023}.\\

\section{Optimización de modelos de \gls{ia} para arquitecturas de recursos limitados}

La viabilidad de ejecutar algoritmos de Inteligencia Artificial para la protección infantil depende directamente de su optimización técnica, dado que los dispositivos móviles presentan severas restricciones térmicas, computacionales y energéticas. Para mitigar el agotamiento prematuro de la batería, es imperativo utilizar arquitecturas ligeras y estrategias de escalado adaptativo. Investigaciones en computación móvil han demostrado que implementar sistemas de decisión dinámica que gestionen la calidad del renderizado visual y la frecuencia de muestreo de los sensores puede lograr ahorros energéticos superiores al 30\% sin sacrificar la precisión de la detección en tiempo real \cite{ali_efficient_2021}. En el desarrollo de nuestra aplicación, la integración de modelos de aprendizaje automático se agiliza utilizando enfoques de Desarrollo Dirigido por Modelos (MDD) y marcos de trabajo de arquitectura limpia (\textit{Clean Architecture}), que permiten generar código optimizado y modular para despliegues móviles multiplataforma \cite{alwakeel_appcraft_2025}.\\

La evolución hacia el procesamiento en el borde, o \textit{Edge Computing}, representa el paradigma definitivo para la \gls{ia} móvil. Al trasladar la inferencia desde servidores remotos hacia el procesador del dispositivo local (\gls{npu} o GPU), se elimina la latencia de red y se garantiza la privacidad de los datos infantiles. Estrategias avanzadas como la cuantización de pesos y la destilación de conocimiento permiten que modelos complejos de visión artificial operen en teléfonos de gama media \cite{surianarayanan_survey_2023}. Este enfoque ha sido validado en dominios computacionalmente exigentes como la Realidad Aumentada Móvil (MAR), donde las interfaces adaptativas al contexto del usuario procesan ricas señales visuales localmente para mantener interacciones fluidas \cite{cao_mobile_2023}. Al aplicar estos mismos principios de optimización visual a la supervisión de pantalla, el sistema logra identificar amenazas sin interrumpir la actividad del menor.\\


Asimismo, la optimización local se beneficia de conceptos extraídos de los sistemas inteligentes distribuidos, como los utilizados en el transporte inteligente y el \gls{iot}, donde los nodos individuales (en este caso, los teléfonos móviles) toman decisiones críticas de forma autónoma sin esperar instrucciones de un servidor central \cite{oladimeji_smart_2023}. Para asegurar que la captura de pantalla y el análisis del texto se realicen en milisegundos, también se aplican técnicas de optimización procedentes del desarrollo de videojuegos en Unity y Android, tales como el agrupamiento de objetos (\textit{object pooling}) y la eliminación de llamadas redundantes a la \gls{api} del sistema operativo, garantizando que el servicio en segundo plano opere con un impacto de \gls{cpu} casi nulo \cite{koulaxidis_improving_2022}. Estas optimizaciones cruzadas son esenciales para que la aplicación no sea desinstalada por causar un deterioro en el rendimiento del dispositivo.\\

\section{Limitaciones técnicas, sesgos y consideraciones éticas del modelo}

A pesar de los sofisticados avances en la detección inteligente, la implementación de \gls{ia} en la supervisión de menores enfrenta importantes limitaciones técnicas y dilemas éticos. El obstáculo algorítmico más frecuente es el desequilibrio en la tasa de falsos positivos y falsos negativos. Una configuración excesivamente sensible puede generar alertas constantes sobre contenidos educativos legítimos o interacciones inofensivas, provocando "fatiga de alerta" en los padres y deteriorando la confianza técnica en el software \cite{rojas_technical_2025, banic_comparison_2024}. Por el contrario, un modelo demasiado permisivo corre el riesgo de omitir amenazas sutiles pero devastadoras, como las micro-agresiones del ciberacoso o la presión de grupo relacionada con el \textit{sexting} y el FOMO (\textit{Fear Of Missing Out}), fenómenos que tienen una alta prevalencia en las interacciones nocturnas de los adolescentes \cite{tomczyk_sexting_2023}.\\

Desde una perspectiva ética, el uso de \gls{ia} predictiva en contextos de protección social y familiar levanta preocupaciones sobre el sesgo algorítmico y la discriminación. Si los modelos de detección son entrenados con conjuntos de datos no representativos, podrían penalizar injustamente ciertos dialectos juveniles o expresiones culturales. Por ello, la literatura académica exige la integración de "justicia algorítmica" y transparencia, requiriendo que los sistemas eviten las decisiones de "caja negra" y expliquen al usuario por qué se desencadenó una alerta \cite{agbana_ethical_2025}. Para que la supervisión móvil sea éticamente aceptable, debe mantener al "humano en el ciclo" (\textit{human-in-the-loop}), permitiendo que la \gls{ia} actúe como un evaluador primario pero cediendo el juicio definitivo y la acción disciplinaria al razonamiento empático del padre.\\


Otra limitación crítica en el despliegue de este modelo es la brecha de usabilidad y alfabetización tecnológica en el entorno familiar. La adopción de comportamientos de ciberseguridad está fuertemente dictada por la percepción del riesgo de privacidad; si los padres no confían en el manejo de datos de la aplicación, rechazarán su uso \cite{balapour_mobile_2020}. Las revisiones de diseño enfatizan que las herramientas de control parental deben integrar recomendaciones claras de usabilidad y seguridad para empoderar a los padres, quienes a menudo carecen de las habilidades para configurar permisos complejos \cite{gnanasekaran_usability_2023}. Estudios recientes confirman que, aunque la inmensa mayoría de los tutores considera la ciberseguridad una prioridad (87\%), la falta de educación y conciencia práctica impide una supervisión efectiva \cite{ayyash_cybersecurity_2024}. Esto se agrava por el hecho de que muchos padres subestiman los rastros digitales que sus hijos dejan en redes sociales y su impacto en vulnerabilidades futuras \cite{chang_cybersecurity_2023}.\\

\section{Intervención pedagógica y mediación habilitadora asistida por \gls{ia}}

El propósito último de integrar Inteligencia Artificial en la supervisión infantil no es consolidar un modelo de vigilancia autoritaria, sino facilitar una "mediación habilitadora" que construya resiliencia. La investigación sociotécnica demuestra que los padres que combinan la restricción técnica con la co-navegación y la comunicación abierta logran maximizar las oportunidades en línea para sus hijos mientras minimizan los daños \cite{livingstone_maximizing_2017}. En este marco, la \gls{ia} asume el rol de un "asistente de mediación", proporcionando a los tutores la evidencia visual necesaria para iniciar diálogos pedagógicos sobre riesgos específicos, un enfoque que ha demostrado ser altamente eficaz en estudios cualitativos de interacción familiar y prevención de la adicción a internet \cite{theopilus_parental_2025, allison_parental_2025}.\\

Esta capacidad de la tecnología para actuar como puente emocional transforma la herramienta en una plataforma de salud mental digital. Durante eventos disruptivos como la pandemia, el uso de dispositivos móviles por parte de los adolescentes funcionó tanto como un mecanismo de aislamiento adictivo como una vía de soporte social; la mediación inteligente ayuda a potenciar los aspectos positivos del uso digital mitigando la exposición a contenidos que inducen estrés o comparaciones nocivas \cite{marciano_digital_2022}. Además, al integrar principios de las intervenciones digitales de salud mental, la aplicación puede fomentar una participación más activa del menor (\textit{engagement}), utilizando la retroalimentación de la \gls{ia} para enseñarles a auto-gestionar sus hábitos y reconocer sus propios límites \cite{liverpool_engaging_2020}.\\


Finalmente, el diseño de soluciones de desconexión tecnológica (\textit{tech disengagement}) para adolescentes tempranos subraya la importancia de evitar medidas exclusivamente restrictivas. Las plataformas que promueven el monitoreo colaborativo y la definición conjunta de reglas fortalecen la capacidad del niño para adaptarse a los entornos digitales novedosos y peligrosos de forma autónoma \cite{chowdhury_landscape_2025, ziker_parentchild_2025}. Al evaluar el bienestar digital subjetivo desde una perspectiva ecológica, se evidencia que los niños requieren que sus propias perspectivas sean consideradas en el diseño de las herramientas que rigen su actividad en línea \cite{holmarsdottir_integrative_2025}. Por tanto, la bitácora de evidencias asistida por \gls{ia} culmina no como un registro policial, sino como un andamiaje educativo que se retira progresivamente a medida que el menor demuestra madurez y ciudadanía digital.